\scanpage{313}
\mspage{263}

\uline{Dém} Considérons une factorisation de $S\longrightarrow X_1$ en

\diagram{\(\begin{tikzcd}[ampersand replacement=\&]
S\arrow[r,"i_1","{\unsure{\in|\mathrm C}}"']\&\widetilde X_1\arrow[r,dash,"p_1"]\&X_1
\end{tikzcd},\)}

avec $i_1$ imm ouverte de Dwyer, $p_1\in W$. Donc \unsure{on a} un carré comm.

\diagram{\(\widetilde Q=\left(\begin{tikzcd}[ampersand replacement=\&,column sep=small,row sep=small]
S\arrow[r,hook,"{\mathrm{Dw}}"]\arrow[d]\&\widetilde X_1\arrow[d]\arrow[r,"{\deleted{\unsure{p_1}}}"]\arrow[r,dash,decorate,decoration={zigzag,segment length=2pt,amplitude=1pt}]\&\deleted{\unsure{X_1}}\arrow[d,"{\deleted{\unsure{1}}}"]\arrow[d,dash,decorate,decoration={zigzag,segment length=3pt,amplitude=2pt}]\\
X_2\arrow[r]\&\deleted{\unsure{X_1}}X_3\arrow[r,"{\deleted{\unsure{\leftarrow}}}"]\&X_3
\end{tikzcd}\right)\)}

et un hom $\widetilde Q\xrightarrow{\alpha}\overset{\deleted{\unsure{\sim}}}{Q}$ qui est l'identité sur \inserted{les} trois \deleted{\unsure{des}} sommets distincts de $\widetilde X_1$, et qui est $p_1$ en $\widetilde X_1$.

\unsure{On} en déduit un diagramme comm.

\diagram{\(\begin{tikzcd}[ampersand replacement=\&,column sep=small,row sep=small]
\unsure{\mathscr X|\mathcal X}_Q\arrow[rr,"{\mathscr X_\alpha=\lambda}"]\arrow[dr,"{\varphi_Q}"']\&{}\&\mathscr X_{\widetilde Q}\arrow[dl,dash,"{\varphi_{\widetilde Q}}"]\\
\&X_3\&{}
\end{tikzcd},\)}

et on a $\lambda\in W$ -- (car \deleted{il \unsure{faut}} c'est un hom \unsure{$\Phi$-fibré|$\Phi$-fibr.} de catégories fibrées sur $\Phi$, qui \deleted{\unsure{induit des}} une W-équi. sur chacune des fibres).

% slanted
\margin{On utilise ici $W(3)$ (forme $\Phi$ forte), i.e. $\mathrm{Loc}(3\ \unsure{bis})$ \unsure{ou} -- \unsure{f}, car \uline{cartésiens}, \unsure{de|ou} \uline{\unsure{XVI}}. Mais \unsure{je pense|j'en passe}, cf plus bas.}

Donc $\varphi_Q\in W\Longleftrightarrow\varphi_{\widetilde Q}$ \unsure{tel}, d'autre part il est \deleted{\unsure{pour prouver}}\inserted{clair que} $Q$ est $W$-cocart. ssi $\widetilde Q$ l'est. Donc \unsure{SPS} que $S\hookrightarrow X_1$ est imm. ouverte (de Dwyer). Appliquant ceci
% archive sheet number: 313
